\section{Setup}
\label{sec:setup}

\paragraph{Model.}
An input $X\in\R^{H\times W\times S}$ has a spectrum $x_{h,w}\in\R^S$ at
every pixel. A \emph{spectral encoder} $f_\thetap:\R^S\to\R^K$ is applied
per pixel, producing $Z\in\R^{K\times H\times W}$; a \emph{spatial model}
$g_\phip$ of characteristic width $M$ (channels, embedding dimension or
hidden width) maps $Z$ to per-pixel logits $\hat y\in\R^{N_{\mathrm{cls}}
\times H\times W}$. We write $F_{\thetap,\phip}=g_\phip\circ f_\thetap$,
$\Cf=\dim\thetap$ and $\Cg=\dim\phip$. The theorems are informative when
$\Cg\gg\Cf$; the production model of Section~\ref{sec:diagnostics} has
$\Cg/\Cf\approx 299$.

\paragraph{Loss, residual and Jacobians.}
With per-pixel cross-entropy averaged over the $N=N_{\mathrm{data}}HW$
pixels, the gradient with respect to the logits at pixel $n$ is the residual
$r_n=\mathrm{softmax}(\hat y_n)-y_{n,\mathrm{onehot}}$. We stack with
$N^{-1/2}$ so that no constant depends on the dataset size:
\begin{equation}
\label{eq:jacobians}
  r:=N^{-1/2}(r_1,\dots,r_N),\qquad
  J_b:=N^{-1/2}\,\frac{\partial(\hat y_1,\dots,\hat y_N)}{\partial b},\qquad
  \nabla_b\loss=J_b^\top r,\quad b\in\{\thetap,\phip\}.
\end{equation}
Two objects must be kept apart. The Gauss--Newton matrix and its diagonal
blocks live in parameter space,
\begin{equation}
\label{eq:ggn}
  G=J^\top H_\tau J,\qquad J=[J_\thetap,\,J_\phip],\qquad
  G_{bb}=J_b^\top H_\tau J_b,
\end{equation}
with $H_\tau=\mathrm{diag}(p)-pp^\top$ the softmax Jacobian, blockwise over
pixels \citep{botev2017practical}; they measure \emph{curvature}. The
output-space kernels $K_b=J_bJ_b^\top$ measure \emph{gradient gain}: for
cross-entropy, $r^\top K_br=\norm{\nabla_b\loss}^2$ holds exactly, with no
squared-loss approximation. Section~\ref{sec:geometry} concerns curvature;
Section~\ref{sec:dynamics} concerns gain along the realised residual.

\begin{definition}[Module curvature disparity]
\label{def:dcurv}
$\Dcurv(G):=\lambda_{\max}(G_{\phip\phip})/\lambda_{\max}(G_{\thetap\thetap})$,
evaluated at random initialization unless stated otherwise, with
$\Dcurv:=+\infty$ if the spectral block vanishes.
\end{definition}

\paragraph{Assumptions.}
\begin{assumption}
\label{ass:linear}
(i) \emph{Linear encoder}: $f_\thetap(x)=W^\top x$ with
$\thetap=\mathrm{vec}(W)$, $W\in\R^{S\times K}$.
\end{assumption}
\begin{assumption}
\label{ass:subg}
(ii) \emph{Bounded input Gram}: the mean per-pixel second moment
$\Sigma_X:=N^{-1}\sum_n x_nx_n^\top$ satisfies
$\lambda_{\max}(\Sigma_X)\le\Lambda_X<\infty$.
\end{assumption}
\begin{assumption}
\label{ass:inputlip}
(iii) \emph{Width-uniform input Jacobian}: there is $\widetilde L$,
independent of $M$ and $\Cg$, with
$\sup\|\partial g_\phip(Z)/\partial Z\|_{\mathrm{op}}\le\widetilde L$ over a
parameter regime $\Phi_{\mathrm{reg}}$ and feature domain
$\mathcal{Z}_{\mathrm{reg}}$ (compact when $g_\phip$ uses attention). This
is a modelling assumption: it is proved in expectation for a two-layer ReLU
head at initialization (Appendix~\ref{app:proofs}), where the naive product
of layer norms would give $\Theta(M/K)$ instead.
\end{assumption}
\begin{assumption}
\label{ass:reg}
(iv) \emph{Regularity}: $g_\phip$ is $C^2$ jointly in $(\phip,Z)$, or
piecewise smooth (ReLU), in which case gradient flow is read almost
everywhere with a consistent gate selection
(Appendix~\ref{app:math}).
\end{assumption}
